% Full solutions to the five VI/17 challenge problems.

\begin{solution}[Challenge VI.17.C1]
Let $D(f_1),\dots,D(f_n)$ cover $X=\Spec A$.  The sheaf-theoretic proof is immediate once the principal-open formula is known: the sheaf axiom for this finite cover says exactly that
\[
A=\Gamma(X,\OO_X)
\longrightarrow
\prod_i A_{f_i}
\rightrightarrows
\prod_{i,j}A_{f_if_j}
\]
is an equalizer.  The two arrows are the two restriction maps, and a tuple lies in the equalizer precisely when its members agree on every pairwise overlap.

For a direct algebraic proof, first prove injectivity.  Suppose $a\in A$ maps to zero in every $A_{f_i}$.  For each $i$ there is $N_i$ with $f_i^{N_i}a=0$.  Choose a common $N$.  Since the opens $D(f_i^N)=D(f_i)$ still cover $X$, the ideal $(f_1^N,\dots,f_n^N)$ is the unit ideal.  Thus there exist $r_i\in A$ with
\[
\sum_i r_i f_i^N=1.
\]
Multiplying by $a$ gives $a=0$.

Now let $(s_i)$ be a compatible tuple with $s_i\in A_{f_i}$.  Because the cover and the set of pairs are finite, after increasing one common exponent $N$ we may write
\[
s_i=\frac{a_i}{f_i^N}
\]
and arrange, by clearing the finitely many overlap equalities, that
\[
f_j^N a_i=f_i^N a_j
\qquad\text{for all }i,j.
\]
Again choose $r_i\in A$ with $\sum_i r_i f_i^N=1$ and set
\[
a=\sum_i r_i a_i.
\]
Fix $j$.  Then
\[
\begin{aligned}
f_j^N a-a_j
&=\sum_i r_i f_j^N a_i-a_j\sum_i r_i f_i^N\\
&=\sum_i r_i\bigl(f_j^N a_i-f_i^N a_j\bigr)=0.
\end{aligned}
\]
Hence $a/1=a_j/f_j^N=s_j$ in $A_{f_j}$.  So $a$ maps to the compatible tuple.  Injectivity gives uniqueness.  This proves the equalizer statement directly in commutative algebra.
\end{solution}

\begin{solution}[Challenge VI.17.C2]
Let
\[
A=k[x_1,\dots,x_n],\qquad n\ge2,
\]
and let $U=\mathbb A_k^n\setminus\{0\}$.  The coordinate principal opens
\[
D(x_1),\dots,D(x_n)
\]
cover $U$: a prime corresponding to a point different from the origin cannot contain all coordinates.

Since $A$ is a domain, all localizations embed in the common fraction field
\[
K=k(x_1,\dots,x_n).
\]
The equalizer description of sections on a principal cover identifies
\[
\Gamma(U,\OO_U)=\bigcap_{i=1}^n A_{x_i}\subseteq K.
\]
We show the intersection is $A$.

Take $r\in\bigcap_iA_{x_i}$ and write $r=p/q$ in lowest terms in the UFD $A$.  Membership $r\in A_{x_i}$ means that after cancellation every irreducible factor of the denominator $q$ is associated to $x_i$; equivalently, $q$ is a unit times a power of $x_i$.  Because this must hold for every $i$, and $n\ge2$, no nonunit irreducible can occur in $q$: the variables $x_1$ and $x_2$ are nonassociate primes.  Hence $q$ is a unit and $r\in A$.  The reverse inclusion $A\subseteq A_{x_i}$ is obvious for every $i$.  Therefore
\[
\Gamma(U,\OO_U)=A.
\]
This is an elementary affine precursor of Hartogs-type extension: removing the codimension-$n$ origin for $n\ge2$ creates no new global regular functions on affine space.
\end{solution}

\begin{solution}[Challenge VI.17.C3]
Assume $I^N=0$ for some $N$.  Every prime ideal of $A$ contains $I$, because every element of $I$ is nilpotent.  Hence the quotient map $A\to A/I$ induces a homeomorphism
\[
\Spec(A/I)\xrightarrow{\sim}\Spec A,
\qquad
\mfp/I\longmapsto\mfp.
\]
We use this homeomorphism to regard the two structure sheaves as sheaves on the same topological space $X$.

At $\mfp\in X$, the map on stalks is
\[
A_{\mfp}\longrightarrow (A/I)_{\mfp/I}.
\]
Localization commutes with quotient, so
\[
(A/I)_{\mfp/I}\cong A_{\mfp}/I_{\mfp}.
\]
Thus the stalk map is surjective and its kernel is exactly $I_{\mfp}$.

Let $\widetilde I$ denote the quasi-coherent ideal sheaf associated with the $A$-module $I$.  Its stalk at $\mfp$ is $I_{\mfp}$.  Since kernels of morphisms of sheaves are determined stalkwise, the sheaf kernel of
\[
\OO_{\Spec A}\longrightarrow\OO_{\Spec(A/I)}
\]
is precisely $\widetilde I$.  Moreover
\[
(\widetilde I)^N=0
\]
stalkwise, so this kernel is a nilpotent ideal sheaf.  Therefore a nilpotent thickening changes no topological points; it enriches their local function rings by a nilpotent sheaf ideal.
\end{solution}

\begin{solution}[Challenge VI.17.C4]
Suppose
\[
X=\Spec A=U\sqcup V
\]
with $U$ and $V$ nonempty and both open and closed.  Define a local section on $U$ to be $1$ and a local section on $V$ to be $0$.  Because $U\cap V=\varnothing$, these sections are automatically compatible.  The sheaf gluing axiom therefore gives a unique global section
\[
e\in\Gamma(X,\OO_X)=A
\]
whose restriction is $1$ on $U$ and $0$ on $V$.

The sections $e^2$ and $e$ agree on $U$ and on $V$, so by locality they are equal globally.  Hence
\[
e^2=e.
\]
Because both $U$ and $V$ are nonempty, $e$ is neither $0$ nor $1$.

We now identify the two opens.  If $\mfp\in U$, then the germ $e_{\mfp}=1$ is a unit in $A_{\mfp}$, so $e\notin\mfp$ and $\mfp\in D(e)$.  If $\mfp\in V$, then $e_{\mfp}=0$.  Were $e\notin\mfp$, the idempotent $e$ would be a unit in the local ring $A_{\mfp}$; a unit cannot be zero.  Thus $e\in\mfp$.  Therefore
\[
U=D(e).
\]
The same argument applied to $1-e$ gives
\[
V=D(1-e).
\]
Thus nontrivial clopen decompositions of an affine spectrum are exactly the geometric form of nontrivial idempotents.
\end{solution}

\begin{solution}[Challenge VI.17.C5]
Let
\[
\Phi:(\Spec A,\OO_A)\xrightarrow{\sim}(\Spec B,\OO_B)
\]
be an isomorphism of locally ringed spaces.  A morphism of ringed spaces contains a morphism of sheaves
\[
\OO_B\longrightarrow \Phi_*\OO_A.
\]
Passing to global sections gives a ring homomorphism
\[
\Gamma(\Spec B,\OO_B)
\longrightarrow
\Gamma(\Spec A,\OO_A).
\]
Because $\Phi$ is an isomorphism, this map is an isomorphism.  Using the affine global-section theorem on both sides,
\[
B\cong\Gamma(\Spec B,\OO_B)
\cong
\Gamma(\Spec A,\OO_A)
\cong A.
\]
Hence $A\cong B$.

This does not yet prove the full anti-equivalence between rings and affine schemes.  We have shown only that an isomorphism of affine locally ringed spaces forces an isomorphism of rings.  Full anti-equivalence requires much more: for every ring homomorphism $B\to A$ one must construct a corresponding morphism
\[
\Spec A\to\Spec B,
\]
prove that its map on structure sheaves is local on stalks, prove uniqueness, and establish a natural bijection
\[
\Hom_{\mathrm{AffSch}}(\Spec A,\Spec B)
\cong
\Hom_{\mathrm{Ring}}(B,A).
\]
That morphism classification belongs to the later affine-scheme chapter rather than to the present structure-sheaf construction.
\end{solution}
